\begin{textsample}{POS Dim 1 – 2007 – Score 32.00 – 2007/t000238.txt}  \label{ex:f1_pos_021}
Well, good morning.
You know, \textbf{the} \textbf{computer} and television both recently turned 60, and today I ' d like to talk about their relationship.
Despite their middle age, if you ' ve been following \textbf{the} \textbf{themes} of this conference or \textbf{the} entertainment industry, it ' s pretty clear that one has been picking on \textbf{the} other.
So it ' s about time that we talked about how \textbf{the} \textbf{computer} ambushed television, or why \textbf{the} invention of \textbf{the} atomic bomb unleashed forces that lead to \textbf{the} writers ' strike.
And it ' s not just what these are doing to each other, but it ' s what \textbf{the} audience thinks that really frames this matter.
To get a sense of this, and it ' s been a \textbf{theme} we ' ve talked about all week, I recently talked to a bunch of tweeners.
I wrote on cards :"television,""radio,""MySpace,""Internet,""PC."And I said, just arrange these, from what ' s important to you and what ' s not, and then tell me why.
Let ' s listen to what happens when they get to \textbf{the} portion of \textbf{the} discussion on television.
Girl 1 : Well, I think it ' s important but, like, not necessary because you can do a lot of other stuff with your free time than watch programs.
Peter Hirshberg : Which is more fun, Internet or TV?
Girls : Internet.
Girl 2 : I think we — \textbf{the} reasons, one of \textbf{the} reasons we put \textbf{computer} before TV is because nowadays, like, we have TV shows on \textbf{the} \textbf{computer}.
Girl 2 : And then you can download onto your iPod.
PH : Would you like to be \textbf{the} president of a TV network?
Girl 4 : I wouldn ' t like it.
Girl 2 : That would be so stressful.
Girl 5 : No.
PH : How come?
Girl 5 : Because they ' re going to lose all their money eventually.
Girl 3 : Like \textbf{the} stock market, it goes up and down and stuff.
I think right now \textbf{the} \textbf{computers} will be at \textbf{the} top and everything will be kind of going down and stuff.
PH : There ' s been an uneasy relationship between \textbf{the} TV business and \textbf{the} tech business, really ever since they both turned about 30.
We go through periods of enthrallment, followed by reactions in boardrooms, in \textbf{the} finance community best characterized as, what ' s \textbf{the} finance term?
Ick pooey.
Let me give you an example of this. \textbf{The} year is 1976, and Warner buys Atari because video games are on \textbf{the} rise. \textbf{The} next year they march forward and they introduce Qube, \textbf{the} first interactive cable TV system, and \textbf{the} New York Times heralds this as telecommunications moving to \textbf{the} home, convergence, great things are happening.
Everybody in \textbf{the} East Coast gets in \textbf{the} pictures — Citicorp, Penney, RCA — all getting into this big vision.
By \textbf{the} way, this is about when I enter \textbf{the} picture.
I ' m going to do a summer internship at Time Warner.
That summer I ' m all — I ' m at Warner that summer — I ' m all excited to work on convergence, and then \textbf{the} bottom falls out.
Doesn ' t work out too well for them, they lose money.
And I had a happy brush with convergence until, kind of, Warner basically has to liquidate \textbf{the} whole thing.
That ' s when I leave graduate school, and I can ' t work in New York on kind of entertainment and technology because I have to be exiled to California, where \textbf{the} remaining jobs are, almost to \textbf{the} \textbf{sea}, to go to work for Apple Computer.
Warner, of course, writes off more than 400 million dollars.
Four hundred million dollars, which was real money back in \textbf{the} ' 70s.
But they were onto something and they got better at it.
By \textbf{the} year 2000, \textbf{the} process was perfected.
They merged with AOL, and in just four years, managed to shed about 200 billion dollars of market capitalization, showing that they ' d actually mastered \textbf{the} art of applying Moore ' s law of successive miniaturization to their balance sheet.
Now, I think that one reason that \textbf{the} media and \textbf{the} entertainment communities, or \textbf{the} media community, is driven so crazy by \textbf{the} tech community is that tech folks talk differently.
You know, for 50 years, we ' ve talked about changing \textbf{the} world, about total transformation.
For 50 years, it ' s been about hopes and fears and promises of a better world.
And I got to thinking, you know, who else talks that way?
And \textbf{the} answer is pretty clearly — it ' s people in religion and in politics.
And so I realized that actually \textbf{the} tech world is best understood, not as a business cycle, but as a messianic movement.
We promise something great, we evangelize it, we ' re going to change \textbf{the} world.
It doesn ' t work out too well, and so we actually go back to \textbf{the} well and start all over again, as \textbf{the} people in New York and L.
A. look on in absolute, morbid astonishment.
But it ' s this irrational view of things that drives us on to \textbf{the} next thing.
So, what I ' d like to ask is, if \textbf{the} \textbf{computer} is becoming a principal tool of media and entertainment, how did we get here?
I mean, how did a machine that was built for accounting and artillery morph into media?
Of course, \textbf{the} first \textbf{computer} was built just after World War II to solve military problems, but things got really interesting just a couple of years later — 1949 with Whirlwind, built at MIT ' s Lincoln Lab.
Jay Forrester was building this for \textbf{the} Navy, but you can ' t help but see that \textbf{the} creator of this machine had in mind a machine that might actually be a potential media star.
So take a look at what happens when \textbf{the} foremost journalist of early television meets one of \textbf{the} foremost \textbf{computer} pioneers, and \textbf{the} \textbf{computer} begins to express itself.
Journalist : It ' s a Whirlwind electronic \textbf{computer}.
With considerable trepidation, we \textbf{undertake} to interview this new machine.
Jay Forrester : Hello New York, this is Cambridge.
And this is \textbf{the} oscilloscope of \textbf{the} Whirlwind electronic \textbf{computer}.
Would you like if I used \textbf{the} machine?
Journalist : Yes, of course.
But I have an idea, Mr.
Forrester.
Since this \textbf{computer} was made in conjunction with \textbf{the} Office of Naval Research, why don ' t we switch down to \textbf{the} Pentagon in Washington and let \textbf{the} Navy ' s research chief, Admiral Bolster, give Whirlwind \textbf{the} workout?
Calvin Bolster : Well, Ed, this problem concerns \textbf{the} Navy ' s Viking rocket.
This rocket goes up 135 miles into \textbf{the} sky.
Now, at \textbf{the} standard rate of fuel \textbf{consumption}, I would like to see \textbf{the} \textbf{computer} trace \textbf{the} flight path of this rocket and see how it can determine, at any instant, say at \textbf{the} end of 40 seconds, \textbf{the} amount of fuel remaining, and \textbf{the} velocity at that set instant.
JF : Over on \textbf{the} left-hand side, you will notice fuel \textbf{consumption} decreasing as \textbf{the} rocket takes off.
And on \textbf{the} \textbf{right-hand} side, there ' s a scale that shows \textbf{the} rocket ' s velocity. \textbf{The} rocket ' s position is shown by \textbf{the} trajectory that we ' re now looking at.
And as it reaches \textbf{the} peak of its trajectory, \textbf{the} velocity, you will notice, has dropped off to a minimum.
Then, as \textbf{the} rocket dives down, velocity picks up again toward a maximum velocity and \textbf{the} rocket hits \textbf{the} ground.
How ' s that?
Journalist : What about that, Admiral?
CB : Looks very good to me.
JF : And before leaving, we would like to show you another kind of mathematical problem that some of \textbf{the} boys have worked out in their spare time, in a less serious vein, for a Sunday afternoon.
Journalist : Thank you very much indeed, Mr.
Forrester and \textbf{the} MIT lab.
PH : You know, so much was worked out : \textbf{the} first real-time interaction, \textbf{the} video display, pointing a gun.
It lead to \textbf{the} microcomputer, but unfortunately, it was too pricey for \textbf{the} Navy, and all of this would have been lost if it weren ' t for a happy coincidence.
Enter \textbf{the} atomic bomb.
We ' re threatened by \textbf{the} greatest weapon ever, and knowing a good thing when it sees it, \textbf{the} Air Force decides it needs \textbf{the} biggest \textbf{computer} ever to protect us.
They adapt Whirlwind to a massive air defense system, deploy it all across \textbf{the} frozen north, and spend nearly three times as much on this \textbf{computer} as was spent on \textbf{the} Manhattan Project building \textbf{the} A-Bomb in \textbf{the} first place.
Talk about a shot in \textbf{the} arm for \textbf{the} \textbf{computer} industry.
And you can imagine that \textbf{the} Air Force became a pretty good salesman.
Here ' s their marketing video.
Narrator : In a mass raid, high-speed bombers could be in on us before we could determine their tracks.
And then it would be too late to act.
We cannot afford to take that chance.
It is to meet this threat that \textbf{the} Air Force has been developing SAGE, \textbf{the} Semi-Automatic Ground Environment system, to strengthen our air defenses.
This new \textbf{computer}, built to become \textbf{the} nerve center of a defense network, is able to perform all \textbf{the} complex mathematical problems involved in countering a mass enemy raid.
It is provided with its own powerhouse containing large diesel-driven generators, air-conditioning equipment, and cooling towers required to cool \textbf{the} thousands of vacuum tubes in \textbf{the} \textbf{computer}.
PH : You know, that one \textbf{computer} was huge.
There ' s an interesting marketing lesson from it, which is basically, when you market a product, you can either say, this is going to be wonderful, it will make you feel better and enliven you.
Or there ' s one other marketing proposition : if you don ' t use our product, you ' ll die.
This is a really good example of that.
This had \textbf{the} first pointing \textbf{device}.
It was distributed, so it worked out — distributed computing and modems — so all these things could talk to each other.
About 20 percent of all \textbf{the} nation ' s programmers were wrapped up in this thing, and it led to an awful lot of what we have today.
It also used vacuum tubes.
You saw how huge it was, and to give you a sense for this — because we ' ve talked a lot about Moore ' s law and making things small at this conference, so let ' s talk about making things large.
If we took Whirlwind and put it in a place that you all know, say, Century City, it would fit beautifully.
You ' d kind of have to take Century City out, but it could fit in there.
But like, let ' s imagine we took \textbf{the} latest Pentium processor, \textbf{the} latest Core 2 Extreme, which is a four-core processor that Intel ' s working on, it will be our \textbf{laptop} tomorrow.
To build that, what we ' d do with Whirlwind technology is we ' d have to take up roughly from \textbf{the} 10 to Mulholland, and from \textbf{the} 405 to La Cienega just with those Whirlwinds.
And then, \textbf{the} 92 nuclear power plants that it would take to provide \textbf{the} power would fill up \textbf{the} rest of Los Angeles.
That ' s roughly a third more nuclear power than all of France creates.
So, \textbf{the} next time they tell you they ' re on to something, clearly they ' re not.
So — and we haven ' t even worked out \textbf{the} cooling needs.
But it gives you \textbf{the} kind of power that people have, that \textbf{the} audience has, and \textbf{the} reasons these transformations are happening.
All of this stuff starts moving into industry.
DEC kind of reduces all this and makes \textbf{the} first mini-computer.
It shows up at places like MIT, and then a mutation happens.
Spacewar! is built, \textbf{the} first \textbf{computer} game, and all of a sudden, interactivity and involvement and passion is worked out.
Actually, many MIT students stayed up all night long working on this thing, and many of \textbf{the} principles of gaming today were worked out.
DEC knew a good thing about wasting time.
It shipped every one of its \textbf{computers} with that game.
Meanwhile, as all of this is happening, by \textbf{the} mid- ' 50s, \textbf{the} business model of traditional broadcasting and cinema has been busted completely.
A new technology has confounded radio men and movie moguls and they ' re quite certain that television is about to do them in.
In fact, despair is in \textbf{the} air.
And a quote that sounds largely reminiscent from everything I ' ve been reading all week.
RCA had David Sarnoff, who basically commercialized radio, said this,"I don ' t say that radio networks must die.
Every effort has been made and will continue to be made to find a new pattern, new selling arrangements and new types of programs that may arrest \textbf{the} declining revenues.
It may yet be possible to eke out a poor existence for radio, but I don ' t know how."And of course, as \textbf{the} \textbf{computer} industry develops interactively, producers in \textbf{the} emerging TV business actually hit on \textbf{the} same idea.
And they fake it.
Jack Berry : Boys and girls, I think you all know how to get your magic windows up on \textbf{the} set, you just get them out.
First of all, get your Winky Dink kits out.
Put out your Magic Window and your erasing glove, and rub it like this.
That ' s \textbf{the} way we get some of \textbf{the} magic into it, boys and girls.
Then take it and put it right up against \textbf{the} screen of your own television set, and rub it out from \textbf{the} center to \textbf{the} corners, like this.
Make sure you keep your magic crayons handy, your Winky Dink crayons and your erasing glove, because you ' ll be using them during \textbf{the} show to draw like that.
You all set?
OK, let ' s get right to \textbf{the} first story about Dusty Man.
Come on into \textbf{the} secret lab.
PH : It was \textbf{the} dawn of interactive TV, and you may have noticed they wanted to sell you \textbf{the} Winky Dink kits.
Those are \textbf{the} Winky Dink crayons.
I know what you ' re saying."Pete, I could use any ordinary \textbf{open-source} crayon, why do I have to buy theirs?"I assure you, that ' s not \textbf{the} case.
Turns out they told us directly that these are \textbf{the} only crayons you should ever use with your Winky Dink Magic Window, other crayons may discolor or hurt \textbf{the} window.
This proprietary principle of vendor lock-in would go on to be perfected with great success as one of \textbf{the} enduring principles of windowing systems everywhere.
It led to lawsuits — — federal investigations, and lots of repercussions, and that ' s a scandal we won ' t discuss today.
But we will discuss this scandal, because this man, Jack Berry, \textbf{the} host of"Winky Dink,"went on to become \textbf{the} host of"Twenty One,"one of \textbf{the} most important quiz shows ever.
And it was rigged, and it became unraveled when this man, Charles van Doren, was outed after an unnatural winning streak, ending Berry ' s career.
And actually, ending \textbf{the} career of a lot of people at CBS.
It turns out there was a lot to learn about how this new medium worked.
And 50 years ago, if you ' d been at a meeting like this and were trying to understand \textbf{the} media, there was one prophet and only but one you wanted to hear from, Professor Marshall McLuhan.
He actually understood something about a \textbf{theme} that we ' ve been discussing all week.
It ' s \textbf{the} role of \textbf{the} audience in an era of pervasive electronic communications.
Here he is talking from \textbf{the} 1960s.
Marshall McLuhan : If \textbf{the} audience can become involved in \textbf{the} actual process of making \textbf{the} ad, then it ' s happy.
It ' s like \textbf{the} old quiz shows.
They were great TV because it gave \textbf{the} audience a role, something to do.
They were horrified when they discovered they ' d really been left out all \textbf{the} time because \textbf{the} shows were rigged.
Now, then, this was a horrible misunderstanding of TV on \textbf{the} part of \textbf{the} programmers.
PH : You know, McLuhan talked about \textbf{the} global village.
If you substitute \textbf{the} word blogosphere, of \textbf{the} Internet today, it is very true that his understanding is probably very enlightening now.
Let ' s listen in to him.
MM : \textbf{The} global village is a world in which you don ' t necessarily have harmony.
You have extreme concern with everybody else ' s business and much involvement in everybody else ' s life.
It ' s a sort of Ann Landers ' column writ large.
And it doesn ' t necessarily mean harmony and peace and quiet, but it does mean huge involvement in everybody else ' s affairs.
And so \textbf{the} global village is as big as a \textbf{planet}, and as small as a village post office.
PH : We ' ll talk a little bit more about him later.
We ' re now right into \textbf{the} 1960s.
It ' s \textbf{the} era of big business and data centers for computing.
But all that was about to change.
You know, \textbf{the} expression of technology reflects \textbf{the} people and \textbf{the} time of \textbf{the} culture it was built in.
And when I say that code expresses our hopes and aspirations, it ' s not just a joke about messianism, it ' s actually what we do.
But for this part of \textbf{the} story, I ' d actually like to throw it to America ' s leading technology correspondent, John Markoff.
John Markoff : Do you want to know what \textbf{the} counterculture in drugs, sex, \textbf{rock} ' n ' roll and \textbf{the} anti-war movement had to do with computing?
Everything.
It all happened within five miles of where I ' m standing, at Stanford University, between 1960 and 1975.
In \textbf{the} midst of revolution in \textbf{the} streets and \textbf{rock} and roll concerts in \textbf{the} parks, a group of researchers led by people like John McCarthy, a \textbf{computer} scientist at \textbf{the} Stanford \textbf{Artificial} Intelligence Lab, and Doug Engelbart, a \textbf{computer} scientist at SRI, changed \textbf{the} world.
Engelbart came out of a pretty dry engineering culture, but while he was beginning to do his work, all of this stuff was bubbling on \textbf{the} mid-peninsula.
There was LSD leaking out of Kesey ' s Veterans ' Hospital experiments and other areas around \textbf{the} campus, and there was music literally in \textbf{the} streets. \textbf{The} Grateful Dead was playing in \textbf{the} pizza parlors.
People were leaving to go back to \textbf{the} land.
There was \textbf{the} Vietnam War.
There was black liberation.
There was women ' s liberation.
This was a remarkable place, at a remarkable time.
And into that ferment came \textbf{the} microprocessor.
I think it was that interaction that led to personal computing.
They saw these tools that were controlled by \textbf{the} establishment as ones that could actually be liberated and put to use by these communities that they were trying to build.
And most importantly, they had this ethos of sharing information.
I think these ideas are difficult to understand, because when you ' re trapped in one paradigm, \textbf{the} next paradigm is always like a science fiction \textbf{universe} — it makes no sense. \textbf{The} stories were so compelling that I decided to write a book about them. \textbf{The} title of \textbf{the} book is,"What \textbf{the} Dormouse Said : How \textbf{the} ' 60s Counterculture Shaped \textbf{the} Personal Computer Industry."\textbf{The} title was taken from \textbf{the} lyrics to a Jefferson \textbf{Airplane} song. \textbf{The} lyrics go,"Remember what \textbf{the} dormouse said.
Feed your head, feed your head, feed your head."PH : By this time, computing had kind of leapt into media territory, and in short order much of what we ' re doing today was imagined in Cambridge and Silicon Valley.
Here ' s \textbf{the} Architecture Machine Group, \textbf{the} predecessor of \textbf{the} Media Lab, in 1981.
Meanwhile, in California, we were trying to commercialize a lot of this stuff.
HyperCard was \textbf{the} first program to introduce \textbf{the} public to hyperlinks, where you could randomly hook to any kind of picture, or piece of text, or data across a file system, and we had no way of explaining it.
There was no metaphor.
Was it a database?
A prototyping tool?
A scripted language?
Heck, it was everything.
So we ended up writing a marketing brochure.
We asked a question about how \textbf{the} mind works, and we let our customers play \textbf{the} role of so many blind men filling out \textbf{the} \textbf{elephant}.
A few years later, we then hit on \textbf{the} idea of explaining to people \textbf{the} secret of, how do you get \textbf{the} content you want, \textbf{the} way you want it and \textbf{the} easy way?
Here ' s \textbf{the} Apple marketing video.
James Burke : You ' ll be pleased to know, I ' m sure, that there are several ways to create a HyperCard interactive video. \textbf{The} most involved method is to go ahead and produce your own videodisc as well as build your own HyperCard \textbf{stacks}.
By far \textbf{the} simplest method is to buy a pre-made videodisc and HyperCard \textbf{stacks} from a commercial supplier. \textbf{The} method we illustrate in this video uses a pre-made videodisc but creates custom HyperCard \textbf{stacks}.
This method allows you to use existing videodisc materials in ways which suit your specific needs and interests.
PH : I hope you realize how subversive that is.
That ' s like a Dick Cheney speech.
You think he ' s a nice balding guy, but he ' s just declared war on \textbf{the} content business.
Find \textbf{the} commercial stuff, mash it up, tell \textbf{the} story your way.
Now, as long as we confine this to \textbf{the} education market, and a personal matter between \textbf{the} \textbf{computer} and \textbf{the} file system, that ' s fine, but as you can see, it was about to leap out and upset Jack Valenti and a lot of other people.
By \textbf{the} way, speaking of \textbf{the} filing system, it never \textbf{occurred} to us that these hyperlinks could go beyond \textbf{the} local area network.
A few years later, Tim Berners-Lee worked that out.
It became a killer app of links, and today, of course, we call that \textbf{the} World Wide Web.
Now, not only was I instrumental in helping Apple miss \textbf{the} Internet, but a couple of years later, I helped Bill Gates do \textbf{the} same thing. \textbf{The} year is 1993 and he was working on a book and I was working on a video to help him kind of explain where we were all heading and how to popularize all this.
We were plenty aware that we were messing with media, and on \textbf{the} \textbf{surface}, it looks like we predicted a lot of \textbf{the} right things, but we also missed an awful lot.
Let ' s take a look.
Narrator : \textbf{The} pyramids, \textbf{the} Colosseum, \textbf{the} New York subway system and TV dinners, ancient and modern wonders of \textbf{the} man-made world all.
Yet each pales to insignificance with \textbf{the} completion of that magnificent accomplishment of twenty-first-century technology, \textbf{the} Digital Superhighway.
Once it was only a dream of technoids and a few long-forgotten politicians. \textbf{The} Digital Highway arrived in America ' s living rooms late in \textbf{the} twentieth century.
Let us recall \textbf{the} pioneers who made this technical marvel possible. \textbf{The} Digital Highway would follow \textbf{the} rutted trail first blazed by Alexander Graham Bell.
Though some were incredulous...
Man 1 : \textbf{The} phone company!
Narrator : Stirred by \textbf{the} prospects of mass communication and making big \textbf{bucks} on advertising, David Sarnoff commercializes radio.
Man 2 : Never had scientists been put under such pressure and demand.
Narrator : \textbf{The} medium introduced America to new products.
Voice 1 : Say, mom, Windows for Radio means more enjoyment and greater ease of use for \textbf{the} whole family.
Be sure to \textbf{enjoy} Windows for Radio at home and at work.
Narrator : In 1939, \textbf{the} Radio Corporation of America introduced television.
Man 2 : Never had scientists been put under such pressure and demand.
Narrator : Eventually, \textbf{the} race to \textbf{the} future took on added momentum with \textbf{the} breakup of \textbf{the} telephone company.
And further stimulus came with \textbf{the} deregulation of \textbf{the} cable television industry, and \textbf{the} re-regulation of \textbf{the} cable television industry.
Ted Turner : We did \textbf{the} work to build this, this cable industry, now \textbf{the} broadcasters want some of our money.
I mean, it ' s ridiculous.
Narrator : \textbf{Computers}, once \textbf{the} unwieldy tools of accountants and other \textbf{geeks}, escaped \textbf{the} backrooms to enter \textbf{the} media fracas. \textbf{The} world and all its culture reduced to bits, \textbf{the} lingua franca of all media.
And \textbf{the} forces of convergence exploded.
Finally, four great industrial sectors combined.
Telecommunications, entertainment, computing and everything else.
Man 3 : We ' ll see channels for \textbf{the} gourmet and we ' ll see channels for \textbf{the} pet lover.
Voice 2 : Next on \textbf{the} gourmet pet channel, decorating birthday cakes for your schnauzer.
Narrator : All of industry was in play, as investors flocked to place their bets.
At stake : \textbf{the} battle for you, \textbf{the} consumer, and \textbf{the} right to spend billions to send a lot of information into \textbf{the} parlors of America.
PH : We missed a lot.
You know, you missed, we missed \textbf{the} Internet, \textbf{the} long tail, \textbf{the} role of \textbf{the} audience, open systems, social networks.
It just goes to show how tough it is to come up with \textbf{the} right uses of media.
Thomas Edison had \textbf{the} same problem.
He wrote a list of what \textbf{the} phonograph might be good for when he invented it, and kind of only one of his ideas turned out to have been \textbf{the} right early idea.
Well, you know where we ' re going on from here.
We come into \textbf{the} era of \textbf{the} dotcom, \textbf{the} World Wide Web, and I don ' t need to tell you about that because we all went through that bubble together.
But when we emerge from this and what we call Web 2. 0, things actually are quite different.
And I think it ' s \textbf{the} reason that TV ' s so challenged.
If Internet one was about pages, now it ' s about people.
It ' s a customer, it ' s an audience, it ' s a person who ' s participating.
It ' s \textbf{the} formidable thing that is changing entertainment now.
MM : Because it gave \textbf{the} audience a role, something to do.
PH : In my own company, Technorati, we see something like 67, 000 blog posts an hour come in.
That ' s about 2, 700 fresh, connective links across about 112 million blogs that are out there.
And it ' s no wonder that as we head into \textbf{the} writers ' strike, odd things happen.
You know, it reminds me of that old saw in Hollywood, that a producer is anyone who knows a writer.
I now think a network boss is anyone who has a cable modem.
But it ' s not a joke.
This is a real headline."Websites attract striking writers : operators of sites like MyDamnChannel. com could benefit from labor disputes."Meanwhile, you have \textbf{the} TV bloggers going out on strike, in sympathy with \textbf{the} television writers.
And then you have TV Guide, a Fox property, which is about to sponsor \textbf{the} online video awards — but cancels it out of sympathy with traditional television, not appearing to gloat.
To show you how schizophrenic this all is, here ' s \textbf{the} head of MySpace, or Fox Interactive, a News Corp company, being asked, well, with \textbf{the} writers ' strike, isn ' t this going to hurt News Corp and help you online?
Man : But I, yeah, I think there ' s an opportunity.
As \textbf{the} strike continues, there ' s an opportunity for more people to experience video on places like MySpace TV.
PH : Oh, but then he remembers he works for Rupert Murdoch.
Man : Yes, well, first, you know, I ' m part of News Corporation as part of Fox Entertainment Group.
Obviously, we hope that \textbf{the} strike is — that \textbf{the} issues are resolved as quickly as possible.
PH : One of \textbf{the} great things that ' s going on here is \textbf{the} globalization of content really is happening.
Here is a \textbf{clip} from a video, from a piece of animation that was written by a writer in Hollywood, animation worked out in Israel, farmed out to Croatia and India, and it ' s now an international series.
Narrator : \textbf{The} following takes place between \textbf{the} minutes of 2:15 p. m. and 2:18 p. m., in \textbf{the} months preceding \textbf{the} presidential primaries.
Voice 1 : You ' ll have to stay here in \textbf{the} safe house until we get word \textbf{the} \textbf{terrorist} threat is over.
Voice 2 : You mean we ' ll have to live here, together?
Voices 2, 3 and 4 : With her?
Voice 2 : Well, there goes \textbf{the} neighborhood.
PH : \textbf{The} company that created this, Aniboom, is an interesting example of where this is headed.
Traditional TV animation costs, say, between 80, 000 and 10, 000 dollars a minute.
They ' re producing things for between 1, 500 and 800 dollars a minute.
And they ' re offering their creators 30 percent of \textbf{the} back end, in a much more entrepreneurial manner.
So, it ' s a different model.
What \textbf{the} entertainment business is struggling with, \textbf{the} world of brands is figuring out.
For example, Nike now understands that Nike Plus is not just a \textbf{device} in its shoe, it ' s a network to hook its customers together.
And \textbf{the} head of marketing at Nike says,"People are coming to our site an average of three times a week.
We don ' t have to go to them."Which means television advertising is down 57 percent for Nike.
Or, as Nike ' s head of marketing says,"We ' re not in \textbf{the} business of keeping media companies alive.
We ' re in \textbf{the} business of connecting with consumers."And media companies realize \textbf{the} audience is important also.
Here ' s a man announcing \textbf{the} new Market Watch from Dow Jones, powered 100 percent by \textbf{the} user experience on \textbf{the} home page — user-generated content married up with traditional content.
It turns out you have a bigger audience and more interest if you hook up with them.
Or, as Geoffrey Moore once told me, it ' s intellectual curiosity that ' s \textbf{the} trade that brands need in \textbf{the} age of \textbf{the} blogosphere.
And I think this is beginning to happen in \textbf{the} entertainment business.
One of my heroes is songwriter, Ally Willis, who just wrote"\textbf{The} Color Purple"and has been an R and — rhythm and \textbf{blues} writer, and this is what she said about where songwriting ' s going.
Ally Willis : Where millions of collaborators wanted \textbf{the} song, because to look at them strictly as spam is missing what this medium is about.
PH : So, to wrap up, I ' d love to throw it back to Marshall McLuhan, who, 40 years ago, was dealing with audiences that were going through just as much change, and I think that, today, traditional Hollywood and \textbf{the} writers are framing this perhaps in \textbf{the} way that it was being framed before.
But I don ' t need to tell you this, let ' s throw it back to him.
Narrator : We are in \textbf{the} middle of a tremendous clash between \textbf{the} old and \textbf{the} new.
MM : \textbf{The} medium does things to people and they are always completely unaware of this.
They don ' t really notice \textbf{the} new medium that is wrapping them up.
They think of \textbf{the} old medium, because \textbf{the} old medium is always \textbf{the} content of \textbf{the} new medium, as movies tend to be \textbf{the} content of TV, and as books used to be \textbf{the} content, novels used to be \textbf{the} content of movies.
And so every time a new medium arrives, \textbf{the} old medium is \textbf{the} content, and it is highly observable, highly noticeable, but \textbf{the} real, real roughing up and massaging is done by \textbf{the} new medium, and it is ignored.
PH : I think it ' s a great time of enthrallment.
There ' s been more raw DNA of communications and media thrown out there.
Content is moving from shows to particles that are batted back and forth, and part of social communications, and I think this is going to be a time of great renaissance and opportunity.
And whereas television may have gotten beat up, what ' s getting built is a really exciting new form of communication, and we kind of have \textbf{the} merger of \textbf{the} two industries and a new way of thinking to look at it.
Thanks very much.

% matched lemmas: airplane, artificial, blue, buck, clip, computer, consumption, device, elephant, enjoy, geek, laptop, occur, open-source, planet, right-hand, rock, sea, stack, surface, terrorist, the, theme, undertake, universe
\end{textsample}
